\chapter{Kategorik Veriler, Çapraz Tablolar ve Ki-Kare Bağımsızlık Testine Giriş}
\label{ch:categorical}

\begin{hedefler}
Bu bölümün sonunda:
\begin{itemize}
  \item çapraz tabloda araştırma sorusuna uygun satır, sütun veya toplam yüzdesini seçebilecek,
  \item bağımsızlık altındaki beklenen frekansları, hücre katkılarını ve artıkları hesaplayabilecek,
  \item uyum iyiliği, bağımsızlık, homojenlik ve Fisher kesin testi arasında seçim yapabilecek,
  \item Pearson ki-kare sonucunu Cramér $V$, hücre örüntüsü ve varsayımlarla raporlayabileceksiniz.
\end{itemize}
\end{hedefler}

\begin{prerequisitebox}
Kategorik değişken, frekans, yüzde, bağımsızlık ve hipotez testi kavramları
bilinmelidir. Satır ve sütun toplamlarıyla oran hesaplanabilmelidir.
\end{prerequisitebox}

\begin{hazirlik}
İki programın başarı yüzdeleri karşılaştırılırken yalnız başarılı öğrenci
sayılarına bakmak neden yanıltıcı olabilir? Grup büyüklüklerini kullanarak
açıklayın.
\end{hazirlik}

Kategorik değişkenlerin ölçme düzeyi Bölüm~\ref{sec:population-data-types}'te,
hipotez testi ve \pval-değeri yorumu Bölüm~\ref{ch:hypothesis}'te verilmiştir.

\section{Kategorik veriyi özetlemek}

Kategorik değişkenlerde temel özet frekans ve yüzdedir. İki kategorik
değişken birlikte incelendiğinde satır ve sütun kategorilerinden oluşan çapraz
tablo kullanılır. Grup büyüklükleri farklıysa yalnız hücre sayılarına değil,
uygun satır veya sütun yüzdelerine bakılmalıdır. Kategorik veri çözümlemesinin
kuramsal ve uygulamalı çerçevesi için bkz. \citet{agresti2019}.

\section{Bağımsızlık hipotezi}

Ki-kare bağımsızlık testinde
\[
H_0:\text{iki kategorik değişken bağımsızdır}
\]
hipotezi sınanır. Bağımsızlık altında $(i,j)$ hücresinin beklenen frekansı
\[
E_{ij}=\frac{R_iC_j}{N}
\]
ile hesaplanır. Burada $R_i$ satır, $C_j$ sütun toplamıdır.

\begin{definitionbox}
\textbf{Ki-kare bağımsızlık testi}, iki kategorik değişkenin ortak
frekanslarını, değişkenler bağımsız olsaydı beklenen frekanslarla karşılaştırır.
Genel test ilişkinin yerini ve büyüklüğünü tek başına göstermez.
\end{definitionbox}

\section{Ki-kare test ailesinde doğru soruyu seçmek}
\label{sec:chi-choice}

\textbf{Uyum iyiliği testi}, tek bir kategorik değişkenin gözlenen
dağılımını önceden belirlenmiş oranlarla karşılaştırır. \textbf{Bağımsızlık
testi}, tek bir evrende ölçülen iki kategorik değişkenin ilişkili olup
olmadığını inceler. \textbf{Homojenlik testi} ise farklı evren veya gruplarda
aynı kategorik değişkenin dağılımlarını karşılaştırır. Bağımsızlık ve
homojenlik testlerinin hesapları aynı çapraz tablo formülünü kullanabilse de
örnekleme tasarımları ve yorumları farklıdır.

Aynı kişilerin önce ve sonra ölçüldüğü ikili kategorik veriler bağımsız
değildir. Böyle bir durumda $2\times2$ bağımlı tablo için McNemar testi gibi
eşleştirilmiş veriye uygun bir yöntem gerekir.

\section{Pearson ki-kare istatistiği}
\label{sec:pearson-chi}

Gözlenen ve beklenen frekanslar arasındaki standartlaştırılmış farklar
\[
\chi^2=\sum_i\sum_j\frac{(O_{ij}-E_{ij})^2}{E_{ij}}
\]
istatistiğinde birleştirilir. $r\times c$ tabloda serbestlik derecesi
$(r-1)(c-1)$'dir.

Büyük $\chi^2$ değeri gözlenen tablonun bağımsızlık altında beklenen tablodan
uzak olduğunu gösterir. Hangi hücrelerin farkı oluşturduğunu anlamak için
standartlaştırılmış artıklar ve tablo yüzdeleri incelenir.

\section{Uygulama koşulları}
\label{sec:chi-assumptions}

Gözlemler bağımsız olmalı ve beklenen frekanslar ki-kare yaklaşımı için
yeterince büyük olmalıdır. Aynı kişinin birden çok kez sayılması bağımsızlığı
bozar. Küçük $2\times2$ tablolarda Fisher kesin testi düşünülebilir
\citep{agresti2018}.

\section{Etki büyüklüğü}
\label{sec:cramers-v}

İlişkinin büyüklüğü Cramér katsayısıyla
\[
V=\sqrt{\frac{\chi^2}{N\min(r-1,c-1)}}
\]
özetlenebilir. Anlamlı test ilişkinin güçlü veya uygulamada önemli olduğunu
tek başına göstermez.

\begin{outputguidebox}
Çıktıda Pearson ki-kare satırını, serbestlik derecesini ve \pval-değerini
okuyun. Beklenen frekans uyarısını, satır/sütun yüzdelerini, artıkları ve
Cramér $V$ değerini ayrıca inceleyin.
\end{outputguidebox}

\section{Raporlama örneği}

``Tercih edilen çözüm stratejisi ile öğretim yöntemi arasında ilişki bulundu,
$\chi^2(2,N=180)=8{,}31$, $p=0{,}016$, $V=0{,}21$. Yüzdeler ve
standartlaştırılmış artıklar farkın özellikle görsel strateji kategorisinde
yoğunlaştığını gösterdi.''

\section{Çözümlü iki çarpı iki tablo}

İki öğretim yönteminde başarı durumu Tablo~\ref{tab:teaching-success}'teki
gibi gözlenmiş olsun.

\begin{table}[H]
\centering
\caption{Öğretim yöntemi ve başarı durumu.}
\label{tab:teaching-success}
\begin{tabular}{lrrr}
\toprule
 & Başarılı & Başarısız & Toplam \\
\midrule
Yöntem A & 30 & 10 & 40 \\
Yöntem B & 20 & 20 & 40 \\
\midrule
Toplam & 50 & 30 & 80 \\
\bottomrule
\end{tabular}
\end{table}

A--başarılı hücresinin beklenen frekansı $40\times50/80=25$'tir. Diğer
beklenen frekanslar 15, 25 ve 15 olur. Pearson testi
$\chi^2(1,N=80)=5{,}33$, $p=0{,}021$ sonucunu verir. Etki büyüklüğü
\[
V=\sqrt{5{,}33/80}=0{,}26
\]
olarak bulunur. A yönteminde başarı oranı yüzde 75, B yönteminde yüzde 50'dir.
İlişki istatistiksel olarak desteklenmektedir; nedensel yorum için yöntemlere
atama sürecinin de incelenmesi gerekir.

\subsection{Yüzdeler ve hücre katkıları}

Araştırma sorusu yöntemler içindeki başarı oranlarını karşılaştırdığı için
satır yüzdeleri kullanılır: A yönteminde $30/40=0{,}75$, B yönteminde
$20/40=0{,}50$. Yalnız ``başarılı öğrencilerin 30'u A yöntemindedir'' demek,
grup büyüklükleri farklı olduğunda yanıltıcı olabilir.

Her hücrenin ki-kare toplamına katkısı $(O-E)^2/E$ ile bulunur ve
Tablo~\ref{tab:chi-cell-contributions}'de ayrı ayrı gösterilir.

\begin{table}[H]
\centering
\caption{Gözlenen, beklenen ve Pearson katkıları.}
\label{tab:chi-cell-contributions}
\begin{tabular}{llrrr}
\toprule
Yöntem & Sonuç & $O$ & $E$ & $(O-E)^2/E$ \\
\midrule
A & Başarılı   & 30 & 25 & 1,00 \\
A & Başarısız  & 10 & 15 & 1,67 \\
B & Başarılı   & 20 & 25 & 1,00 \\
B & Başarısız  & 20 & 15 & 1,67 \\
\midrule
 & Toplam & & & 5,33 \\
\bottomrule
\end{tabular}
\end{table}

Şekil~\ref{fig:chi-residuals}'in temelini oluşturan Pearson artığı
\[
r_{ij}=\frac{O_{ij}-E_{ij}}{\sqrt{E_{ij}}}
\]
ile tanımlanır. Bu değerler sırasıyla yaklaşık $1{,}00$, $-1{,}29$,
$-1{,}00$ ve $1{,}29$'dur. Bu örnekte hiçbir tek hücre olağanüstü büyük bir
artık üretmemekte; genel fark dört hücreye dağılmaktadır. Daha büyük tablolarda
artıklar hangi kategori birleşimlerinin genel ilişkiye yön verdiğini anlamada
özellikle önemlidir.

Pearson artıkları genel tanı için yararlıdır. Satır ve sütun oranlarını da
dikkate alan düzeltilmiş standartlaştırılmış artık
\[
r_{ij}^{\ast}=\frac{O_{ij}-E_{ij}}
{\sqrt{E_{ij}(1-R_i/N)(1-C_j/N)}}
\]
ile hesaplanabilir. Mutlak değeri yaklaşık 2'yi aşan artıklar, genel ilişkinin
hangi hücrelerde yoğunlaşabileceğine işaret eder. Çok sayıda hücre ayrı ayrı
incelendiğinde yanlış pozitif yorum riski arttığından artıklar mekanik birer
hipotez testi gibi kullanılmamalıdır.

\begin{figure}[H]
\centering
\begin{tikzpicture}[x=1.55cm,y=1.0cm]
  \draw[->,PrismGray] (0.2,0)--(5.2,0);
  \draw[->,PrismGray] (0.5,-2.4)--(0.5,2.5)
    node[above,font=\scriptsize]{Pearson artığı};
  \draw[PrismGray,dashed] (0.5,2)--(5.0,2);
  \draw[PrismGray,dashed] (0.5,-2)--(5.0,-2);
  \fill[PrismBlue!78] (0.8,0) rectangle (1.5,1.00);
  \fill[PrismWarnOrange!88!black] (1.9,0) rectangle (2.6,-1.29);
  \fill[PrismWarnOrange!88!black] (3.0,0) rectangle (3.7,-1.00);
  \fill[PrismBlue!78] (4.1,0) rectangle (4.8,1.29);
  \node[below,align=center,font=\tiny] at (1.15,-0.08) {A\\başarılı};
  \node[above,align=center,font=\tiny] at (2.25,0.08) {A\\başarısız};
  \node[above,align=center,font=\tiny] at (3.35,0.08) {B\\başarılı};
  \node[below,align=center,font=\tiny] at (4.45,-0.08) {B\\başarısız};
  \node[left,font=\tiny] at (0.45,2) {$+2$};
  \node[left,font=\tiny] at (0.45,-2) {$-2$};
\end{tikzpicture}
\caption{Çözümlü çapraz tablodaki hücrelerin Pearson artıkları.}
\label{fig:chi-residuals}
\end{figure}

\Needspace{10\baselineskip}
\begin{assumptionbox}
Ki-kare testi frekanslara uygulanır; yüzde veya ortalamalar doğrudan hücre
sayısı gibi kullanılamaz. Her gözlem yalnız bir hücreye katkı vermeli ve
beklenen frekansların dağılımı yaklaşımı desteklemelidir. Aynı öğrencinin iki
zamandaki kategorik yanıtları bağımsızlık testi için bağımsız gözlemler değildir.
\end{assumptionbox}

\begin{mistakebox}
Anlamlı ki-kare sonucu ilişkinin yönünü tek başına göstermez. Hangi grupta
hangi kategorinin daha sık olduğunu anlamak için satır veya sütun yüzdeleri
ile hücre artıklarının incelenmesi gerekir.
\end{mistakebox}

\Needspace{26\baselineskip}
\begin{softwarebox}
SPSS'te \emph{Analyze $>$ Descriptive Statistics $>$ Crosstabs}; Python'da
\texttt{scipy.stats.chi2\_contingency}; R'de \texttt{chisq.test} kullanılabilir.
Yates düzeltmesi gibi yazılım varsayılanları nedeniyle küçük sayısal farklar
oluşabileceğinden kullanılan seçenek raporda belirtilmelidir.

\begin{lstlisting}
import numpy as np
from scipy import stats

observed = np.array([[30, 10],
                     [20, 20]])
chi2, p, dof, expected = stats.chi2_contingency(
    observed, correction=False)
contribution = (observed - expected)**2 / expected
residual = (observed - expected) / np.sqrt(expected)
cramers_v = np.sqrt(chi2 / observed.sum())

print(chi2, p, dof, cramers_v)
print(expected)
print(contribution, residual)
\end{lstlisting}
\end{softwarebox}

\section{Sonuçtan araştırma kararına}

Anlamlı bir ki-kare sonucu iki değişken arasında evrensel ve değişmez bir
ilişki bulunduğunu göstermez. Örnekleme yöntemi, eksik gözlemler, kategori
tanımları ve olası karıştırıcı değişkenler değerlendirilmelidir. Ayrıca
ilişkinin eğitimsel önemini anlayabilmek için grup yüzdeleri ve $V$ değeri
test sonucuyla birlikte sunulmalıdır.

\section{Çözümlü problemler}

\subsection*{Problem 1: Araştırma sorusuna uygun testi seçme}

Aşağıdaki araştırma durumlarının her biri için uygun yöntemi belirleyin.

\begin{enumerate}[label=\alph*)]
  \item Öğrencilerin seçtiği üç öğrenme biçiminin oranları yüzde 50, yüzde 30
  ve yüzde 20 biçimindeki kuramsal dağılımla karşılaştırılıyor.
  \item Okul türü ile tercih edilen öğrenme biçiminin ilişkisi tek bir öğrenci
  örnekleminde inceleniyor.
  \item Üç farklı üniversitedeki öğrencilerin aynı memnuniyet kategorilerine
  dağılımları karşılaştırılıyor.
  \item Aynı öğrencilerin bir eğitimden önce ve sonra başarılı/başarısız
  durumları karşılaştırılıyor.
\end{enumerate}


\textbf{Çözüm.} İlk durumda tek kategorik değişken gözlenen oranlarla
karşılaştırıldığı için ki-kare uyum iyiliği testi uygundur. İkinci durumda
aynı örneklemdeki iki kategorik değişken için ki-kare bağımsızlık testi,
üçüncü durumda farklı grupların dağılımlarını karşılaştırmak için ki-kare
homojenlik testi kullanılır. Dördüncü durumda ölçümler aynı öğrencilere ait
olduğundan bağımsızlık koşulu sağlanmaz; ikili sonuçlar için McNemar testi
düşünülmelidir.

\textbf{Sık hata.} Bir tablonun $2\times2$ olması tek başına Pearson ki-kare
testini uygun kılmaz. Gözlemlerin bağımsız mı, eşleştirilmiş mi olduğu veri
toplama tasarımından belirlenir.

\subsection*{Problem 2: $2\times2$ tabloda tam Pearson çözümü}

Tablo~\ref{tab:teaching-success}'teki Yöntem A ve Yöntem B verilerini
kullanarak bütün beklenen
frekansları, Pearson ki-kare değerini, serbestlik derecesini ve Cramér $V$
değerini hesaplayın.

\textbf{Çözüm.} Satır toplamları 40 ve 40, sütun toplamları 50 ve 30,
genel toplam 80'dir. Beklenen tablo
\[
\begin{pmatrix}
40(50)/80 & 40(30)/80\\
40(50)/80 & 40(30)/80
\end{pmatrix}
=
\begin{pmatrix}25&15\\25&15\end{pmatrix}
\]
olur. Hücre katkılarının toplamı
\[
\chi^2=\frac{(30-25)^2}{25}+\frac{(10-15)^2}{15}
+\frac{(20-25)^2}{25}+\frac{(20-15)^2}{15}=5{,}33
\]
ve serbestlik derecesi $(2-1)(2-1)=1$'dir. Çift yönlü Pearson sonucu
$p\approx0{,}021$ verir. Etki büyüklüğü
\[
V=\sqrt{\frac{5{,}33}{80}}\approx0{,}26
\]
olarak bulunur. Yöntem içi başarı yüzdeleri sırasıyla yüzde 75 ve yüzde
50'dir. Test ilişkiye ilişkin kanıt verir; ilişkinin yönü ve uygulamalı anlamı
yüzdelerden okunmalıdır.

\subsection*{Problem 3: Üç çarpı iki tabloda artık analizi}

Üç öğretim yaklaşımında öğrencilerin görsel veya sözel strateji tercihleri
aşağıdaki gibi gözlenmiştir:
\[
\begin{array}{c|rr|r}
 & \text{Görsel} & \text{Sözel} & \text{Toplam}\\ \hline
A&40&20&60\\
B&30&30&60\\
C&20&40&60\\ \hline
\text{Toplam}&90&90&180
\end{array}
\]
Beklenen frekansları, test sonucunu, Cramér $V$ değerini ve hücre örüntüsünü
yorumlayın.

\textbf{Çözüm.} Her satır toplamı 60 ve her sütun toplamı 90 olduğundan bütün
hücrelerin beklenen frekansı $60(90)/180=30$'dur. A ve C satırlarındaki dört
dış hücrenin her biri $100/30=3{,}33$, B satırındaki hücreler ise sıfır katkı
verir. Böylece
\[
\chi^2=13{,}33,\qquad sd=(3-1)(2-1)=2,
\qquad p\approx0{,}0013.
\]
Cramér katsayısı
\[
V=\sqrt{\frac{13{,}33}{180\min(2,1)}}\approx0{,}27
\]
olur. Pearson artıkları A--görsel ve C--sözel hücrelerinde yaklaşık $+1{,}83$,
A--sözel ve C--görsel hücrelerinde $-1{,}83$'tür. Düzeltilmiş artıkların
mutlak değeri bu dört hücrede yaklaşık 3,16'dır. A grubunda görsel, C grubunda
sözel tercihin bağımsızlık altında beklenenden fazla olduğu; B grubunun ise
beklenen dağılıma uyduğu görülür.

\subsection*{Problem 4: Küçük beklenen frekansta Fisher kesin testi}

Yirmi öğrencilik küçük bir çalışmada sonuçlar
\[
\begin{array}{c|rr}
 & \text{Başarılı} & \text{Başarısız}\\ \hline
A&1&9\\
B&7&3
\end{array}
\]
biçimindedir. Pearson yaklaşımının koşullarını ve uygun analizi değerlendirin.

\textbf{Çözüm.} Bağımsızlık altında beklenen tablo
\[
\begin{pmatrix}4&6\\4&6\end{pmatrix}
\]
olur. Dört hücrenin ikisinde beklenen frekans 5'in altındadır. Bu küçük
$2\times2$ tabloda asimptotik Pearson yaklaşımına güvenmek yerine Fisher kesin
testi tercih edilir. İki yönlü Fisher testi yaklaşık $p=0{,}020$ verir.
Örneklemde başarı oranları A grubunda yüzde 10, B grubunda yüzde 70'tir.

Fisher testi küçük hücre sorununu ele alır; gözlemlerin bağımsız olmaması,
yanlı örnekleme veya gruplara sistematik atama gibi tasarım sorunlarını
düzeltmez. Ayrıca odds oranının hangi grup ve sonuç yönünde tanımlandığı açıkça
belirtilmelidir.

\subsection*{Problem 5: Uyum iyiliği testi}

Yüz yirmi öğrencinin tercih ettiği öğrenme biçiminin eşit oranlarda olması
beklenmektedir. Gözlenen frekanslar yüz yüze için 54, karma için 42 ve çevrim
içi için 24'tür. Eşit dağılım hipotezini sınayın ve etkinin kaynağını açıklayın.

\textbf{Çözüm.} Üç kategori eşit bekleniyorsa her kategorinin beklenen
frekansı $120/3=40$'tır. Test istatistiği
\[
\chi^2=\frac{(54-40)^2}{40}+\frac{(42-40)^2}{40}
+\frac{(24-40)^2}{40}=4{,}90+0{,}10+6{,}40=11{,}40
\]
ve serbestlik derecesi $3-1=2$ olur. \pval-değeri yaklaşık 0,0033'tür; eşit
dağılım hipotezi reddedilir. Pearson artıkları yüz yüze, karma ve çevrim içi
için sırasıyla yaklaşık 2,21, 0,32 ve $-2{,}53$'tür. Fark özellikle yüz yüze
tercihin beklenenden fazla, çevrim içi tercihin beklenenden az olmasından
kaynaklanmaktadır.

Burada beklenen oranların veri görüldükten sonra eşit seçilmemiş olması;
kuram, önceki bilgi veya araştırma tasarımıyla önceden gerekçelendirilmesi
gerekir.

\subsection*{Problem 6: Yazılım çıktısını bütüncül yorumlama}

Bir yazılım çıktısında $3\times2$ tablo için
$\chi^2(2,N=240)=7{,}68$, $p=0{,}021$ ve $V=0{,}18$ verilmiştir. Hücrelerin
yüzde 16,7'sinde beklenen frekans 5'in altında, en küçük beklenen frekans
3,8'dir. Düzeltilmiş artıkların en büyüğü 2,7'dir. Sonucu değerlendirin.

\textbf{Çözüm.} Genel test, iki kategorik değişkenin bağımsız olmadığına
ilişkin yüzde 5 düzeyinde kanıt sunmaktadır. $V=0{,}18$ ilişkinin büyüklüğünü
özetler; ancak küçük, orta veya büyük gibi etiketler araştırma bağlamından
bağımsız kullanılmamalıdır. En büyük mutlak düzeltilmiş artık 2,7 olduğundan
genel ilişkinin belirli bir hücrede yoğunlaştığı düşünülebilir; yönü söylemek
için artığın işareti ve ilgili satır/sütun yüzdesi görülmelidir.

Beklenen frekans özeti tamamen göz ardı edilmemelidir. Hücrelerin yüzde
16,7'sinin 5'in altında olması yaygın öğretici kuraldaki yüzde 20 sınırının
altında kalsa da en küçük değerin 3,8 olması tablo yapısı ve örnekleme planıyla
birlikte değerlendirilmelidir. Gerekirse kategorilerin bilimsel anlamı
bozulmadan birleştirilmesi, kesin veya Monte Carlo \pval-değeri gibi duyarlılık
analizleri raporlanabilir. Yalnız \pval-değerine dayanarak ilişkinin güçlü ya da
nedensel olduğu söylenemez.

\begin{outputguidebox}
Çapraz tablo çıktısını okurken önce gözlenen frekanslarla satır veya sütun
yüzdelerini, sonra beklenen frekans uyarılarını inceleyin. Pearson değerini,
serbestlik derecesini ve \pval-değerini Cramér $V$ ile birlikte raporlayın. Genel
test anlamlıysa ilişkinin nerede yoğunlaştığını artıklar üzerinden, yönünü ise
yüzdeler üzerinden açıklayın.
\end{outputguidebox}

\section{Literatür verisiyle yazılım uygulamaları}
\label{sec:spss-b12}

\textbf{Amaç ve veri.} \texttt{b12.csv} ile kaynak cinsiyet kategorisi
ve yıl sonu notunun en az 10 olması arasındaki ilişki incelenir
\citep{cortez2008data}. Satırlar F ve M, sütunlar 10 puanın altı ve
10 ve üzeridir. Sıfır hipotezi bu iki kategorik değişkenin bağımsız
olmasıdır; alternatif hipotez ilişki bulunmasıdır. Anlamlılık düzeyi
0,05'tir. Önceden seçilen yöntem, süreklilik düzeltmesi uygulanmayan
Pearson ki-kare testidir; B11'deki ortalamalar farkı testi değildir.

\textbf{Ortak veri.} Karşılaştırmada kullanılan dosya
\nolinkurl{companion/spss/csv/b12.csv} olup 395 satır ve 10 sütun
içerir. \texttt{sex} için 1=F ve 2=M; \texttt{pass10} için
0: $\texttt{g3}<10$, 1: $\texttt{g3}\geq10$ tanımlıdır.
Sıfır notlu 38 kayıt eksik sayılmaz; eşik altı kategorisinde tutulur.
B06'nın seçim göstergeleri ve ağırlığı kullanılmaz. Paylaşılan CSV
kopyası proje verisiyle aynıdır; bu kontrol Excel içe aktarmanın
doğrulandığı anlamına gelmez. Kodları \texttt{companion} klasörünü
içeren paket kökünden çalıştırın; veri sözlüğü için bkz.
sayfa~\pageref{sec:spss-guide}.

\subsection{IBM SPSS uygulaması}
\textbf{Başlamadan önce.} Önce aşağıdaki veri açma bloğunu
\texttt{EXECUTE} satırı dahil çalıştırın. B11'den kalan etkin veri
yerine B12 dosyası yeniden açılır. Filtre, ağırlık ve dosya bölme
kapalı olmalıdır. Dosya açılamazsa \texttt{/FILE} yolunu düzeltip
bu bloğu yeniden çalıştırın; eski etkin veriyle devam etmeyin.

\begin{spssadimlar}
\spssadim{\emph{Analyze $\to$ Descriptive Statistics $\to$ Crosstabs} yolunu açın.}
\spssadim{\texttt{sex} değişkenini \emph{Row(s)}, \texttt{pass10} değişkenini \emph{Column(s)} alanına taşıyın.}
\spssadim{\emph{Statistics} içinde \emph{Chi-square} ve \emph{Phi and Cramér's V} seçeneklerini işaretleyip \emph{Continue} seçin.}
\spssadim{\emph{Cells} içinde \emph{Observed}, \emph{Expected}, \emph{Row} yüzdeleri ve \emph{Standardized} artıkları seçin. Düzeltilmiş standartlaştırılmış artıklar bu karşılaştırmanın ölçüsü değildir.}
\spssadim{\emph{Continue $\to$ OK} ile çıktıyı alın. \emph{Case Processing Summary} tablosunda 395 geçerli ve sıfır eksik kayıt bulunduğunu kontrol edin.}
\spssadim{\emph{Pearson Chi-Square} satırındaki \emph{Asymptotic Significance (2-sided)} sütununu ve beklenen frekans dipnotunu okuyun. \emph{Symmetric Measures} tablosundan Cramér $V$ değerini alın.}
\end{spssadimlar}

{\raggedright
\noindent\textbf{Komutların görevi.} \texttt{TABLES=sex BY pass10}
satır ve sütunu, \texttt{CHISQ PHI} test ve ilişki büyüklüklerini
belirler. \texttt{CELLS=COUNT EXPECTED ROW SRESID}, gözlenen ve
beklenen sayımları, satır yüzdelerini ve Pearson hücre artıklarını
ister.\par}

\par\noindent\begin{minipage}{\linewidth}
\textbf{Uygulama kodu:} \texttt{b12}. Veri açma ve tanımlama:

\begin{lstlisting}[style=prismSPSS]
SET DECIMAL=DOT.
GET DATA /TYPE=TXT
 /FILE='companion/spss/csv/b12.csv'
 /ENCODING='UTF8'
 /ARRANGEMENT=DELIMITED /DELCASE=LINE
 /DELIMITERS="," /QUALIFIER='"'
 /FIRSTCASE=2
 /VARIABLES=id F8.0 school F8.0 sex F8.0 age F8.0
 studytime F8.0 absences F8.0 g1 F8.0 g2 F8.0
 g3 F8.0 pass10 F8.0.
FILTER OFF.
WEIGHT OFF.
SPLIT FILE OFF.
VARIABLE LABELS
 sex 'Kaynak verideki cinsiyet kategorisi'
 pass10 'Yil sonu notu en az 10'.
VALUE LABELS
 sex 1 'F' 2 'M'
 /pass10 0 '10 puanin alti' 1 '10 ve uzeri'.
VARIABLE LEVEL sex pass10 (NOMINAL).
EXECUTE.
\end{lstlisting}
\end{minipage}\par

\par\noindent\begin{minipage}{\linewidth}
\textbf{Frekanslar ve ilişki testi.} Veri açma bloğundan sonra çalıştırın.

\begin{lstlisting}[style=prismSPSS]
FREQUENCIES VARIABLES=sex pass10.
CROSSTABS /TABLES=sex BY pass10
 /FORMAT=AVALUE TABLES
 /STATISTICS=CHISQ PHI
 /CELLS=COUNT EXPECTED ROW SRESID.
\end{lstlisting}
\end{minipage}\par

\begin{outputguidebox}
\textbf{Paylaşılan SPSS tablolarıyla doğrulanan değerler.}
\emph{Case Processing Summary} tablosunda 395 geçerli ve sıfır
eksik kayıt vardır. Çapraz tabloda F satırı $(75,133)$, M satırı
$(55,132)$; satır toplamları 208 ve 187'dir. Eşiği karşılayanların
satır yüzdeleri sırasıyla yüzde 63,9 ve 70,6 görünür.

\emph{Pearson Chi-Square} satırında $\chi^2(1)=1{,}970$ ve
$p=0{,}160$; dipnotta en küçük beklenen frekans 61,54 ve beklenen
frekansı 5'in altında olan hücre sayısı sıfırdır. \emph{Symmetric
Measures} tablosunda Phi ve Cramér $V$ üç basamakla 0,071'dir.

\emph{Continuity Correction} satırındaki $p=0{,}195$ ve Fisher
testinin iki yönlü $p=0{,}165$ değeri farklı yöntemlere aittir;
Pearson sonucunun yazılımlar arasında uyuşmadığı anlamına gelmez.
Bu karşılaştırmada yöntem, sonuçlara göre değiştirilmez.
\end{outputguidebox}

\par\noindent\begin{minipage}{\linewidth}
\subsection{R uygulaması}
\label{sec:r-gercek-b12}

\texttt{correct=FALSE} süreklilik düzeltmesini kapatır. Aşağıdaki
blokları sırayla ve çok satırlı komutları bütünüyle çalıştırın.

\begin{lstlisting}[style=prismR]
veri <- read.csv("companion/spss/csv/b12.csv")
stopifnot(nrow(veri) == 395, ncol(veri) == 10,
  !anyNA(veri), !anyDuplicated(veri$id),
  all(veri$sex %in% c(1, 2)),
  all(veri$pass10 %in% c(0, 1)),
  all(veri$pass10 == as.integer(veri$g3 >= 10)))
tablo <- table(
  sex=factor(veri$sex, levels=c(1, 2),
             labels=c("F", "M")),
  pass10=factor(veri$pass10, levels=c(0, 1),
                labels=c("10_alti", "10_ve_uzeri")))
sonuc <- chisq.test(tablo, correct=FALSE)
print(c(satir=nrow(veri), sutun=ncol(veri),
  eksik=sum(is.na(veri)), sifir_not=sum(veri$g3 == 0)))
print(addmargins(tablo))
\end{lstlisting}
\end{minipage}\par

\textbf{Dosya yolu sorunu yaşarsanız.} İlk satırın yerine aşağıdaki
komutu çalıştırıp \texttt{b12.csv} dosyasını seçin; sonra kontrol
bloğundan devam edin.
\begin{lstlisting}[style=prismR]
veri <- read.csv(file.choose())
\end{lstlisting}
Ham CSV dosyası analiz çıktısı değildir. Özellikle
\texttt{), digits=15)} gibi kapanış satırlarını tek başına çalıştırmayın;
ilgili \texttt{print} komutunun tamamını seçin.

\par\noindent\begin{minipage}{\linewidth}
\textbf{Karşılaştırılacak sayısal çıktılar.} Satır yüzdelerinin
paydaları F için 208, M için 187'dir. \texttt{residuals}, Pearson
hücre artıklarıdır; R'deki \texttt{stdres} ile aynı değildir.

\begin{lstlisting}[style=prismR]
print(100 * prop.table(tablo, margin=1), digits=15)
print(sonuc$expected, digits=15)
print(sonuc$residuals, digits=15)
print(c(
  n=sum(tablo),
  ki_kare=unname(sonuc$statistic),
  df=unname(sonuc$parameter),
  p=sonuc$p.value,
  min_beklenen=min(sonuc$expected),
  bes_alti_hucre=sum(sonuc$expected < 5),
  cramer_v=sqrt(unname(sonuc$statistic)/sum(tablo))
), digits=15)
\end{lstlisting}

\textbf{Çıktıyı okuma.} Paylaşılan R metninde dört satır yüzdesi,
dört beklenen sayım, dört Pearson artığı ve son özetin yedi değeri
yer alır. $\chi^2=1{,}969808570687404$, $p=0{,}160468182571215$
ve $V=0{,}070617682919937$ bulunmuştur. R burada yeniden
çalıştırılmamış; kullanıcının paylaştığı çıktı karşılaştırılmıştır.
\end{minipage}\par

\par\noindent\begin{minipage}{\linewidth}
\subsection{Python uygulaması}
\label{sec:python-b12}

\texttt{correction=False} aynı Pearson testini seçer. Satır ve
sütun sıraları açıkça sabitlenir. Blokları sırayla çalıştırın.

\begin{lstlisting}[style=prismPythonApp]
import pandas as pd
import numpy as np
from scipy import stats

veri = pd.read_csv("companion/spss/csv/b12.csv")
assert veri.shape == (395, 10)
assert not veri.isna().any().any()
assert veri.id.is_unique
assert veri.sex.isin([1, 2]).all()
assert veri.pass10.isin([0, 1]).all()
assert veri.pass10.eq((veri.g3 >= 10).astype(int)).all()
tablo = pd.crosstab(veri.sex, veri.pass10).reindex(
    index=[1, 2], columns=[0, 1], fill_value=0)
tablo.index = ["F", "M"]
tablo.columns = ["10_alti", "10_ve_uzeri"]
ki_kare, p_degeri, serbestlik, beklenen = (
    stats.chi2_contingency(tablo, correction=False))
yuzdeler = 100 * tablo.div(tablo.sum(axis=1), axis=0)
artiklar = (tablo.to_numpy()-beklenen)/np.sqrt(beklenen)
print("satir:", len(veri), "sutun:", veri.shape[1])
print("eksik:", int(veri.isna().sum().sum()))
print("sifir_not:", int((veri.g3 == 0).sum()))
print(tablo)
print("Satir toplamlari:", tablo.sum(axis=1).to_dict())
print("Sutun toplamlari:", tablo.sum(axis=0).to_dict())
\end{lstlisting}
\end{minipage}\par

\par\noindent\begin{minipage}{\linewidth}
\textbf{Ayrıntılı çıktı.} Aşağıdaki yazdırma düzeni, paylaşılan
Python terminal çıktısındaki 13 ondalık basamaklı değerleri verir.

\begin{lstlisting}[style=prismPythonApp]
for baslik, degerler in [
    ("Satir yuzdeleri", yuzdeler.to_numpy()),
    ("Beklenen sayimlar", beklenen),
    ("Pearson hucre artiklari", artiklar)
]:
    print(baslik)
    cerceve = pd.DataFrame(degerler,
        index=tablo.index, columns=tablo.columns)
    print(cerceve.to_string(
        float_format=lambda sayi: f"{sayi:.13f}"))

ozet = dict(n=int(tablo.to_numpy().sum()),
    ki_kare=ki_kare, df=serbestlik, p=p_degeri,
    min_beklenen=beklenen.min(),
    bes_alti_hucre=int((beklenen < 5).sum()),
    cramer_v=np.sqrt(ki_kare/len(veri)))
for ad, deger in ozet.items():
    print(f"{ad}: {deger:.13f}")
\end{lstlisting}

\textbf{Çıktıyı okuma.} Paylaşılan Python ekranında 395 satır,
10 sütun, sıfır eksik değer ve 38 sıfır not vardır. Sayımlar ve
toplamlar SPSS ile aynıdır. $p=0{,}1604681825712$ ve
$V=0{,}0706176829199$ değerlerini ilgili etiketlerle okuyun.
Kod ham kayıtları değiştirmez; tablo ayrı bir nesnedir.
\end{minipage}\par

\Needspace{12\baselineskip}
\subsection{Sonuçların karşılaştırılması ve ortak rapor}
13 Eylül 2026'da paylaşılan SPSS tabloları, R metni ve Python
terminal ekranı karşılaştırılmıştır. R metnindeki 19 sayısal
özet, Python sonuçlarıyla $10^{-12}$ mutlak tolerans içinde
uyumludur. Python kodu proje CSV'siyle ayrıca çalıştırılmıştır;
ekrandaki 13 ondalık basamaklı sonuçlarla eşleşir. SPSS tabloları
gösterilen yuvarlama hassasiyetlerinde aynı sonuçları verir.
Aşağıdaki tablolar çıktıların kitap düzenine uyarlanmış özetidir;
ekran görüntüsü değildir. Ara hesaplar yuvarlanmamış değerlerle yapılır.

\begin{table}[H]
\centering
\caption{B12'de kaynak kategorisine göre 10 puan eşiği}
\label{tab:b12-dogrulanmis-sayim}
\begin{tabular}{@{}lrrrr@{}}
\toprule
Grup & $<10$ & $\geq10$ & Toplam & $\geq10$ yüzdesi \\
\midrule
F & 75 & 133 & 208 & 63,942308 \\
M & 55 & 132 & 187 & 70,588235 \\
\midrule
Toplam & 130 & 265 & 395 & 67,088608 \\
\bottomrule
\end{tabular}
\par\smallskip
{\footnotesize Son sütun ilgili satırın yüzdesidir: F için
$100(133/208)$, M için $100(132/187)$, toplam için $100(265/395)$.
Tam 10 puan, eşiği karşılayan kategoriye dahildir.\par}
\end{table}

\begin{table}[H]
\centering
\caption{B12'de beklenen sayımlar ve Pearson hücre artıkları}
\label{tab:b12-dogrulanmis-artik}
\begin{tabular}{@{}lrr@{}}
\toprule
Hücre & Beklenen sayım ($E$) & Pearson artığı \\
\midrule
F, $<10$ & 68,455696 & $0{,}790968$ \\
F, $\geq10$ & 139,544304 & $-0{,}553997$ \\
M, $<10$ & 61,544304 & $-0{,}834199$ \\
M, $\geq10$ & 125,455696 & $0{,}584276$ \\
\bottomrule
\end{tabular}
\par\smallskip
{\footnotesize $E=(\text{satır toplamı})(\text{sütun toplamı})/395$;
Pearson artığı $(O-E)/\sqrt{E}$'dir. SPSS \texttt{SRESID}, R
\texttt{residuals} ve Python \texttt{artiklar} burada aynı ölçüdür.
Yuvarlanmamış dört artığın kareleri toplamı Pearson ki-kare değerini verir.\par}
\end{table}

\begin{table}[H]
\centering
\caption{B12'de Pearson testi ve ilişki büyüklüğü}
\label{tab:b12-dogrulanmis-test}
\begin{tabular}{@{}lr@{}}
\toprule
Ölçü & Değer \\
\midrule
Geçerli kayıt sayısı & 395 \\
Pearson $\chi^2$ & 1,969809 \\
Serbestlik derecesi & 1 \\
$p$ değeri & 0,160468 \\
En küçük beklenen frekans & 61,544304 \\
Beklenen frekansı 5'in altındaki hücre & 0 \\
Cramér $V$ & 0,070618 \\
\bottomrule
\end{tabular}
\par\smallskip
{\footnotesize Süreklilik düzeltmesi uygulanmamıştır. Bu $2\times2$
tabloda Cramér $V=\sqrt{\chi^2/395}$'tir; yön belirtmez.\par}
\end{table}

\Needspace{12\baselineskip}
\begin{outputguidebox}
Tablo~\ref{tab:b12-dogrulanmis-sayim}, eşiği karşılayanların
örneklemdeki oranlarını betimler. Yüzdeler farklı olsa da
Tablo~\ref{tab:b12-dogrulanmis-test} içindeki
$p=0{,}160468>0{,}05$ nedeniyle bağımsızlık sıfır hipotezi
reddedilmez. Bu, bağımsızlığın kanıtı veya oranların kesin eşitliği
değildir. Cramér $V=0{,}070618$, ilişki büyüklüğünün örneklem
özetidir; nedensel etki ya da açıklanan varyans yüzdesi değildir.

Tablo~\ref{tab:b12-dogrulanmis-artik} içindeki pozitif artık,
bağımsızlık altında beklenenden fazla; negatif artık daha az
gözlem olduğunu gösterir. Bunlar düzeltilmiş artıklar değildir;
her hücreyi ayrı anlamlılık testi olarak raporlamayın.

Beklenen sayımların yeterli olması, öğrencilerin bağımsızlığını
veya iki okulla sınırlı verinin temsil gücünü doğrulamaz. Test okul,
önceki başarı ve diğer karıştırıcılar için düzeltme içermez.
B11'de yıl sonu \emph{ortalamaları}, burada ise 10 puan eşiğini
karşılama \emph{oranları} karşılaştırılır. Farklı sorulara farklı
kararlar çıkması çelişki değildir; notları iki kategoriye ayırmak
puan farklılıklarına ilişkin bilgiyi kaybettirir.
\end{outputguidebox}

\Needspace{10\baselineskip}
\textbf{Raporlama örneği (doğrulanmış çıktılarla).}
``395 kayıtta yıl sonu notu en az 10 olanların oranı F grubunda
$133/208$ (yüzde 63,942308), M grubunda $132/187$ (yüzde 70,588235)
bulunmuştur. Süreklilik düzeltmesiz Pearson testinde kaynak cinsiyet
kategorisi ile eşiği karşılama durumu arasında yüzde 5 düzeyinde
yeterli ilişki kanıtı elde edilmemiştir
($\chi^2(1)=1{,}969809$, $p=0{,}160468$, Cramér $V=0{,}070618$).
En küçük beklenen frekans 61,544304'tür; 5'in altında beklenen
frekans yoktur. Sıfır notlu 38 kayıt analizde tutulmuştur.
Sonuç bağımsızlığı kanıtlamaz. Gözlemlerin bağımsızlığı varsayılmış;
okul ve diğer karıştırıcılar için düzeltme yapılmamıştır. İki okulla
sınırlı veriden nedensel veya geniş bir evrene yönelik sonuç çıkarılmamıştır.''

\Needspace{8\baselineskip}
\subsection{Karşılaştırmalı görev: aynı öğrenciler, farklı hedefler}
\label{sec:b12-karsilastirmali-gorev}

\textbf{Hedef ve teslim.} B11'in yalnız F--M Welch analiziyle B12'nin
Pearson analizini karşılaştırın; B11'deki eşleştirilmiş dönem farkını bu
iki soruya karıştırmayın. Mevcut çıktıları kullanarak aşağıdaki beş
maddeyi yanıtlayın. Yeni eşik veya yeni anlamlılık testi seçmeyin.
Kontrol yanıtını okumadan önce kendi gerekçeli raporunuzu tamamlayın.

\begin{enumerate}
  \item \textbf{Ne değişti?} İki analizde kayıtlar, gruplar, yanıt
  değişkeni, hedeflenen büyüklük ve sıfır hipotezini karşılaştırın.
  G3'ü iki kategoriye ayırmak hangi bilgiyi kaybettirir?
  \item \textbf{Payda ve yön.} F ve M için eşiği karşılama yüzdelerini
  payları ve paydalarıyla gösterin. F eksi M oran farkını yüzde puan
  olarak hesaplayın. Bu fark B11'in not puanı farkıyla aynı şey midir?
  Cramér $V$ bu farkın işaretini verir mi?
  \item \textbf{Hesap kontrolü.} Yuvarlanmamış Pearson hücre
  artıklarının karelerini toplayıp $\chi^2$ ile karşılaştırın.
  En küçük beklenen sayımı ve 5'in altındaki hücre sayısını belirtin.
  Bu kontrol öğrencilerin bağımsızlığını da doğrular mı?
  \item \textbf{Karşılaştırmalı rapor.} İki soruyu, test kararlarını
  ve ilişki/fark özetlerini içeren kısa bir paragraf yazın.
  ``Welch anlamlı, ki-kare anlamsız; yazılımlar çelişiyor ve ikinci
  test grupların eşitliğini kanıtlıyor'' yorumunu düzeltin.
  \item \textbf{Analiz disiplini.} ``Eşiği değiştirip anlamlı sonuç
  bulalım; sıfır notları da çıkaralım'' önerisini değerlendirin.
  Aynı öğrencilerle yapılan iki analizin neden bağımsız iki tekrar
  çalışma olmadığını ve nedensel sonuç vermediğini açıklayın.
\end{enumerate}

\Needspace{8\baselineskip}
\textbf{Gerekçeli kontrol yanıtı.}
\begin{enumerate}
  \item Her iki analiz de aynı 395 öğrenciyi, 208 F ve 187 M kaydını
  ve sıfır notları kullanır. Welch sayısal G3 için
  $H_0:\mu_{G3,F}=\mu_{G3,M}$ hipotezini sınar. Pearson analizinde
  yanıt $\mathbb{1}(G3\geq10)$'dur; sıfır hipotezi kaynak kategorisi
  ile eşik durumunun bağımsızlığıdır. Bu iki gruplu tabloda bu,
  grupların eşiği karşılama olasılıklarının eşitliğine karşılık gelir.
  Aynı eşik kategorisindeki farklı notlar artık ayırt edilmez.
  Değişen kayıtlar değil, soru, yanıtın gösterimi ve yöntemdir.
  \item Yüzdeler $100(133/208)=63{,}942308$ ve
  $100(132/187)=70{,}588235$'tir. F eksi M farkı yaklaşık
  $-6{,}645928$ yüzde puandır; iki payda da 395 değildir.
  Bu betimsel oran farkı, B11'in $-0{,}948092$ not puanı farkından
  farklı bir büyüklüktür. Cramér $V=0{,}070618$ yön belirtmez.
  Burada oran farkı için ayrıca güven aralığı hesaplanmamıştır.
  \item Dört artığın kareleri toplamı
  $\sum (O-E)^2/E=1{,}969808570687404$ olup Pearson istatistiğidir.
  En küçük beklenen sayım yaklaşık 61,544304, 5'in altındaki hücre
  sayısı sıfırdır. Hücre hesaplarının uygunluğu öğrencilerin
  bağımsızlığını veya örneklemin temsil gücünü kanıtlamaz.
  Yuvarlanmış artıklarla oluşan küçük fark, yöntem çelişkisi değildir.
  \item ``Yıl sonu ortalamaları için F eksi M farkı
  $-0{,}948092$ puandır (\%95 GA
  $[-1{,}850732;\ -0{,}045452]$); Welch testinde
  $p=0{,}039577$ ile eşit ortalamalar hipotezi reddedilmiştir.
  Ayrı eşik sorusunda oranlar yüzde 63,942308 ve 70,588235'tir;
  Pearson testinde $\chi^2(1)=1{,}969809$, $p=0{,}160468$ ve
  $V=0{,}070618$ bulunmuş, bağımsızlık hipotezi reddedilmemiştir.
  Kararlar her test için ayrı 0,05 düzeyindedir; sorular farklı
  olduğundan çelişki yoktur. İkinci sonuç eşitliğin veya bağımsızlığın
  kanıtı değildir.'' Bu karar ayrılığı tek başına yöntemlerden
  birinin üstünlüğünü de göstermez.
  \item Sonuca bakarak eşik seçmek önceden tanımlanan soruyu
  değiştirir; yalnız anlamlı sonucu sunmak uygun değildir.
  Sıfırları dışlamak kayıt kapsamını ve hedefi değiştirir; sıfır
  notlar kendiliğinden eksik değer sayılmaz. Aynı öğrencilerin
  yeniden analizi yeni bağımsız kanıt üretmez. Okul, önceki başarı
  ve diğer karıştırıcılar denetlenmediğinden nedensel açıklama
  yapılamaz; iki okulla sınırlı kaynağın temsil gücü ayrıca tartışılır.
\end{enumerate}

\Needspace{8\baselineskip}
\textbf{Değerlendirme ölçütleri (10 puan).} Her görev maddesindeki
iki bileşenin her biri 1 puandır; doğru ve gerekçeli bileşen 1, eksik
veya yanlış bileşen 0 puan alır. Eşdeğer doğru ifadeler ve belirtilen
basamaklara uygun yuvarlama kabul edilir.
\begin{enumerate}
  \item Ortak kayıt/grupları ve değişen yanıtı ayırma (1);
  iki sıfır hipotezini ve kategorileştirmenin bilgi kaybını açıklama (1).
  \item Pay, payda, yüzdeler ve yüzde puan farkını doğru verme (1);
  not puanı farkından ayrımı ve $V$'nin yönsüzlüğünü açıklama (1).
  \item Artık kareleri toplamını ve beklenen sayım kontrolünü
  doğru verme (1); bunları bağımsızlık kanıtından ayırma (1).
  \item İki sorunun sayısal özetleri ve kararlarını doğru
  raporlama (1); çelişki ve eşitlik kanıtı yorumlarını düzeltme (1).
  \item Eşik değiştirme ve sıfırları silme önerilerinin soru/kapsam
  üzerindeki etkisini açıklama (1); bağımsız tekrar ve nedensellik
  iddialarını sınırlılıklarla birlikte reddetme (1).
\end{enumerate}
Yalnız $p$ değerlerini doğru kopyalamak yeterli değildir. Bağımsızlığı
kanıtlanmış sayan veya nedensel iddiayı koruyan yanıt tam puan alamaz.
Bu ölçütler öğretim amaçlıdır; öğrenme etkileri henüz öğrenci
uygulamasıyla değerlendirilmemiştir.

\section{R ile çözümlü uygulama}
\label{sec:r-b12}

\textbf{Soru.} Birinci grupta 30 başarılı, 10 başarısız; ikinci grupta
20 başarılı, 20 başarısız gözlem vardır. Pearson bağımsızlık testini,
beklenen frekansları ve Cramér $V$ değerini bulun.

\begin{lstlisting}[style=prismR]
gozlenen <- matrix(c(30, 10, 20, 20), nrow = 2,
                    byrow = TRUE,
                    dimnames = list(
                      grup = c("Birinci", "Ikinci"),
                      sonuc = c("Basarili", "Basarisiz")))
sonuc <- chisq.test(gozlenen, correct = FALSE)
print(sonuc)
sonuc$expected
katkilar <- (gozlenen - sonuc$expected)^2 / sonuc$expected
katkilar
sonuc$residuals
cramer_v <- sqrt(as.numeric(sonuc$statistic) /
                 (sum(gozlenen) * (min(dim(gozlenen)) - 1)))
cramer_v
prop.table(gozlenen, margin = 1)
stopifnot(all(sonuc$expected >= 5))
\end{lstlisting}

\begin{outputguidebox}
\textbf{Çözüm ve kontrol değerleri.} Her satır için beklenen sayılar
$(25,15)$; hücre katkıları $(1;1{,}6667)$'dir. Böylece
$\chi^2(1)=5{,}3333$, $p=0{,}02092$ ve $V=0{,}2582$ olur. Başarı oranları
birinci grupta 0,75, ikinci grupta 0,50'dir. Pearson artıkları ilk satırda
$(1;-1{,}2910)$, ikinci satırda $(-1;1{,}2910)$'dır.
\end{outputguidebox}

\textbf{Yorum.} \texttt{correct = FALSE} bölümdeki düzeltmesiz Pearson
hesabını üretir; ikiye iki tablolarda varsayılan süreklilik düzeltmesiyle aynı
sonuç beklenmez. Bağımsız gözlemler ve uygun örnekleme modeli gerekir;
beklenen frekans denetimi yalnız sayısal koşullardan biridir.

\textbf{Pekiştirme ve yanıt.} Sonuç nesnesinin \texttt{residuals} bileşeni
Pearson artıklarını, \texttt{stdres} bileşeni farklı bir ölçeklendirmeyi verir. Bunları aynı
artık türü gibi raporlamayın. Ki-kare testi ilişkinin nedensel olduğunu göstermez.

\begin{etkinlik}
Bir ki-kare sonucunu kapatmadan önce kontrol edin: Doğru yüzde yönü seçildi mi?
Beklenen frekanslar yeterli mi? Farkı hangi hücreler taşıyor? Cramér $V$ ve
tasarım sınırı raporlandı mı?
\end{etkinlik}

\section{Bölüm özeti}

\begin{itemize}
  \item Çapraz tablo frekanslarla kurulur; yorum araştırma sorusuna uygun yüzdelerle yapılır.
  \item Beklenen frekanslar satır ve sütun toplamlarından bağımsızlık altında hesaplanır.
  \item Pearson istatistiği hücre katkılarının toplamıdır; artıklar farkın kaynağını gösterir.
  \item Cramér $V$, ilişkinin büyüklüğünü örneklem hacminden ayrı değerlendirmeye yardım eder.
  \item Uyum iyiliği, bağımsızlık ve homojenlik testleri benzer hesaplar
  kullansa da farklı araştırma sorularını yanıtlar.
  \item Küçük $2\times2$ tablolarda Fisher kesin testi, eşleştirilmiş ikili
  verilerde ise McNemar testi gibi tasarıma uygun yöntemler düşünülmelidir.
\end{itemize}

\bolumkaynaklari{\cite{agresti2018,agresti2019}}

\section*{Alıştırmalar}

\begin{enumerate}
  \item Bir çapraz tabloda neden yüzdelerin de incelenmesi gerektiğini açıklayın.
  \item Beklenen frekans formülünü kullanarak bir hücre için hesaplama yapın.
  \item Anlamlı genel testten sonra artıkların rolünü belirtin.
  \item Çözümlü tabloda Yöntem A--başarısız hücresinin beklenen değerini hesaplayın.
  \item Ki-kare sonucunun neden tek başına nedensel etki göstermediğini açıklayın.
  \item Çözümlü örnekte dört hücre katkısını toplayarak $\chi^2$ değerini doğrulayın.
  \item Aynı öğrencilerin dönem başı ve dönem sonu tercihlerini incelemek için bağımsızlık varsayımını tartışın.
  \item Anlamlı ki-kare sonucunda Cramér $V$ değerinin neden gerekli olduğunu açıklayın.
  \item Uyum iyiliği, bağımsızlık ve homojenlik testlerine birer araştırma
  sorusu örneği verin.
  \item Satır toplamı 50, sütun toplamı 80 ve genel toplamı 200 olan bir
  hücrenin bağımsızlık altındaki beklenen frekansını bulun.
  \item $3\times4$ çapraz tablonun serbestlik derecesini hesaplayın.
  \item $\chi^2=12$, $N=300$ ve $2\times3$ tablo için Cramér $V$ değerini bulun.
  \item Beklenen frekansları 3, 7, 12 ve 18 olan $2\times2$ tabloda Pearson
  yaklaşımının uygunluğunu tartışın.
  \item Aynı öğrencilerin dönem başı ve sonundaki başarılı/başarısız
  durumları için neden McNemar testinin düşünülmesi gerektiğini açıklayın.
  \item Mutlak düzeltilmiş artığı 3,1 olan bir hücrenin ne gösterdiğini ve ne
  göstermediğini belirtin.
  \item Anlamlı ki-kare sonucu ve $V=0{,}06$ bulunan çok büyük bir örneklemi
  istatistiksel ve uygulamalı önem bakımından yorumlayın.
\end{enumerate}

\bolumbaglantisi{chapters/12-kategorik-kikare.tex}